\documentclass[11pt]{article}

\usepackage[margin=1in]{geometry}
\usepackage{graphicx}
\usepackage{booktabs}
\usepackage{tabularx}
\usepackage{amsmath}
\usepackage[round]{natbib}
\usepackage{hyperref}
\hypersetup{hidelinks}  % no colored boxes around refs/cites (cleaner in print)
% --- global preamble hygiene (URL/path breaking, margin overflow) ---
\usepackage{xurl}
\usepackage[htt]{hyphenat}
\sloppy
\emergencystretch=3em
\usepackage{placeins}   % \FloatBarrier keeps figures from drifting past their section
% Give LaTeX more latitude so large floats land near their text instead of the end.
\renewcommand{\topfraction}{0.9}
\renewcommand{\bottomfraction}{0.85}
\renewcommand{\textfraction}{0.08}
\renewcommand{\floatpagefraction}{0.75}

% All numbers/tables below are generated by report/build_report.py from the
% pipeline outputs -- nothing is hand-entered. Re-run that script to refresh.
\input{generated/macros.tex}

% Include a generated map if make_maps.py has produced it, else a placeholder
% box so the document still compiles.
\newcommand{\mapfigure}[3]{%
  \begin{figure}[htbp]\centering
  \IfFileExists{figures/#1}%
    {\includegraphics[width=\linewidth]{figures/#1}}%
    {\fbox{\parbox[c][4.5cm][c]{0.92\linewidth}{\centering\itshape
      [Figure pending] #2\\[2pt]\normalfont\small
      generate with \texttt{build\_figures.py} (writes the PNG to
      \texttt{report/figures/})}}}%
  \caption{#2}\label{#3}\end{figure}}

% Plain figure with caller-supplied \includegraphics options (for the
% non-full-width district bar charts). #1 file, #2 opts, #3 caption, #4 label.
\newcommand{\figwide}[4]{%
  \begin{figure}[htbp]\centering
  \IfFileExists{figures/#1}%
    {\includegraphics[#2]{figures/#1}}%
    {\fbox{\parbox[c][3cm][c]{0.9\linewidth}{\centering\itshape [Figure pending] #3}}}%
  \caption{#3}\label{#4}\end{figure}}

\title{A Report on Allegheny County Property Reassessments\\
       \large Updated automated-valuation analysis with equity testing}
\author{Russell Richie\thanks{Progress and Poverty Institute. Updates and extends the
  June 2026 analysis by Pro-Housing Pittsburgh's Jack Billings and Connor Schwartz.}}
% NOTE: adjust authorship as appropriate.
\date{\today}

\begin{document}
\maketitle

\begin{abstract}
\noindent Allegheny County has gone nearly half a century without regular
countywide property reassessment, leaving assessed values badly detached from
market value. Using the open-source OpenAVMKit mass-appraisal toolkit and
publicly available data, we build automated valuation models for the county's
\totalParcels\ parcels across \numModelGroups\ model groups and evaluate them
with IAAO-standard ratio studies. Current assessments sit at roughly
\countyUnderPct\% of estimated market value (a countywide new-to-current ratio
of \countyUnderRatio). Our models bring the median sales ratio close to~1.0 and
materially improve uniformity. We additionally test the equity of the resulting
assessments across neighborhood income and racial composition: assessment
\emph{levels} are equitable on both axes, but assessment \emph{uniformity} is
weaker in lower-income and majority-minority neighborhoods --- a data-driven
limitation we document rather than a bias in the model inputs.
\end{abstract}

\section{Executive summary}
Property taxes in Allegheny County are levied on assessed land and building
values that, absent regular reassessment, drift ever further from market value.
We use recent arm's-length sales to fit mass-appraisal models for every parcel
in the county and compare the resulting values against both recent sales and the
county's current assessments.

Two findings stand out. First, the current assessment roll is deeply and
unevenly stale: across model groups the county's median sales ratio runs from
\msrCountyUrban\ to \msrCountyCommercial, i.e.\ properties are assessed at well
under their market value, and a revaluation would raise total residential
assessed value by roughly a factor of \countyUnderRatio. Second, our models are
accurate enough to be useful---median sales ratios near 1.0 and, for most
residential groups, coefficients of dispersion (COD) within or close to IAAO
tolerances---while still leaving clear room for improvement, particularly for
commercial property and in thinly traded neighborhoods.

Two practical messages follow for policymakers. A reassessment would shift tax
burden \emph{within} the county---broadly, toward the historically under-assessed
City of Pittsburgh and older river and mill towns, and away from the affluent
outer suburbs---but because Pennsylvania's anti-windfall law forces offsetting cuts
in millage rates, most individual tax bills would change only modestly. And the
work is affordable: building on existing county data and open-source tools, we
estimate a recurring countywide reassessment can be run for on the order of
\$1\,million per year, far below the tens of millions quoted by commercial vendors
(Section~\ref{sec:cost}).

\section{Background}
Allegheny County last conducted a countywide reassessment in 2012, and before that
in 2001; until the 1970s it had reassessed roughly every three years
\citep{briem2022,hopton2018}. Pennsylvania is unusual in not mandating periodic
reassessment, and because peer jurisdictions such as Philadelphia do reassess on a
schedule, Allegheny is among the largest jurisdictions in the United States that
does not \citep{pa_cca,pa_sccal}. Statutory authority for countywide assessment
rests with the county itself rather than its municipalities or school districts.

Left unchecked, the gap between assessed and market value widens every year. That
drift imposes real and well-documented costs: it shifts tax burden arbitrarily
between owners of otherwise-similar property, distorts investment and mobility
decisions, fuels a steady stream of assessment appeals, and erodes confidence in
the system---costs that ultimately make housing more expensive \citep{horton2024}.
Pittsburgh's abandonment of its split-rate (land-heavy) tax in 2001 compounded the
problem by leaning more heavily on a drifting, flat-rate property base
\citep{hopton2018,hughes2006}.

These problems have prompted sustained calls for reform from journalists, civic
groups, and elected officials \citep{wesa2023,prohousing2024,pgh_council2025,postgazette2025}.
In March 2026, County Council members Dan Grzybek, Lissa Geiger Shulman, and
Bethany Hallam introduced Bill 13892-26, which would require regular countywide
reassessments every three years beginning in 2028. This report demonstrates how
open-source automated valuation methods can be applied to existing county data to
(i)~produce a preliminary reassessment, (ii)~quantify how far the current roll has
drifted, (iii)~surface the equity questions any reassessment program must answer,
and (iv)~estimate what such a program would cost.

\section{Methodology}
The International Association of Assessing Officers (IAAO) recognizes three
approaches to valuation: the cost, sales-comparison, and income approaches
\citep{iaao_mass}. Our models rely primarily on the sales approach, with limited
use of the income approach for commercial parcels.

\subsection{Total property valuation}
We build our models with OpenAVMKit \citep{openavmkit}, an open-source
mass-appraisal toolkit. For each of the \numModelGroups\ model groups we train
several engines---multiple linear-regression variants, a distance-weighted
comparable-sales engine, and gradient-boosted decision trees---and combine them
into an ensemble. Models are trained on valid (open-market, arm's-length) sales
and target a common valuation date of \valuationDate; the ratio studies below
use sales over the \studyPeriod\ study period.

Residential models use parcel and building characteristics (land area, finished
living area, building age, quality and condition grades, a condition/desirability
grade, and parcel geometry), location (latitude/longitude, polar coordinates,
neighborhood, census tract, school district, and municipality), a spatial lag of
recent nearby sale prices, proximity to rivers, water bodies, and parks, and
terrain slope. Census tract-level neighborhood context---median household income
and median gross rent---and a market-tier indicator are also included; their use
is paired with the equity testing reported in Section~\ref{sec:equity}. The
commercial model adds commercial rent, employment density, and building
footprint.

\paragraph{Ensemble and vertical equity.}
OpenAVMKit's default ensemble is tuned for accuracy alone, which we found leaves
the blend regressive (cheaper properties over-assessed relative to expensive
ones). Restricting the residential ensemble to the median of the log-scale
regression engines and the strongest tree engine removes most of that
regressivity at no cost to dispersion. We report the resulting vertical-equity
index (VEI; 0 is neutral, negative is regressive) alongside the accuracy metrics.

\subsection{Land valuation}
To support land-value analysis we decompose each parcel's value into land and
improvement components using a neighborhood-level land-share method
\citep{doucet_land}: within each census tract we apply a land-to-total ratio drawn
from existing assessments to the modeled total value, cap land value at total
value, and treat the remainder as improvement value.

\subsection{Data sources}
Parcel characteristics, assessments, and sales come from the county assessment
roll via the Western Pennsylvania Regional Data Center \citep{wprdc}. Census tract
geography and the American Community Survey supply neighborhood context;
OpenStreetMap supplies amenity and water proximity; building footprints and
employment density come from public datasets \citep{sardari2025,census_lodes}; and
commercial rents are gathered from public listings \citep{loopnet}. Input data are
not redistributed in this repository.

\subsection{Measuring assessment quality}
Fair assessments should be accurate (median sales ratio near 1.0), vertically
equitable (parcels assessed proportionally across price levels), and
horizontally equitable (similar parcels assessed similarly). The IAAO summary
statistic for horizontal uniformity is the coefficient of dispersion (COD; lower
is more uniform), for which the IAAO benchmarks are roughly 15 or below for
single-family residential property and up to about 20 for small income-producing
(e.g.\ multi-family) residential property \citep{iaao_ratio}. For vertical equity
we report a vertical-equity index (VEI; 0 neutral, negative regressive) rather
than the IAAO price-related differential. Section~\ref{sec:results} reports the
median sales ratio (MSR), COD, and VEI; Section~\ref{sec:equity} stratifies them
by neighborhood income and racial composition.

\section{Results}\label{sec:results}
Table~\ref{tab:results} reports, for each model group, the median sales ratio and
trimmed COD of both the current county assessments and our new assessments, plus
the vertical-equity index of the new residential models. Ratios are computed
against raw (non-time-adjusted) sale prices; because the models target the
\valuationDate\ valuation date while sales span the prior year, a well-calibrated
model can sit slightly below~1.0.

\begin{table}[htbp]\centering
\caption{Median sales ratio, COD, and VEI: current vs.\ new assessments.}
\label{tab:results}
\input{generated/results_table.tex}
\end{table}

The new assessments move the median sales ratio from far below~1.0 to near~1.0
for every residential group, and improve uniformity for all of them except the
pre-war group, whose COD is essentially unchanged. The commercial group remains
poorly estimated: with only \salesCommercial\ study-period sales it is data-
starved, and we recommend its values be treated as indicative only. Taken
together with the countywide new-to-current ratio of \countyUnderRatio, the
results confirm both that the current roll is severely under-assessed and that
the new assessments are substantially more accurate and (for residential) more
uniform.

\section{Discussion}

\subsection{Neighborhood-level sales ratio analysis}
A central problem with the current roll is inconsistency across areas: many
neighborhoods are under-assessed relative to others. Mapping the median sales
ratio by census tract shows this dispersion under the current assessments
(Figure~\ref{fig:sr-existing}) and its marked reduction under the new
assessments (Figure~\ref{fig:sr-new}).

\mapfigure{median_sales_ratio_existing.png}{Median sales ratio by census tract, current assessments}{fig:sr-existing}
\mapfigure{median_sales_ratio_new.png}{Median sales ratio by census tract, new assessments}{fig:sr-new}

\subsection{Land value analysis}
Aggregated to the census-tract level, current land assessments already show the
expected gradient---highest downtown and in the East End, lower in outlying
areas (Figure~\ref{fig:land-existing}). The new model re-estimates land value per
square foot directly from land and vacant-parcel sales. Because such sales are
sparse, it currently publishes a tract-level land value only where enough of them
exist to identify one, so its map covers a minority of tracts
(Figure~\ref{fig:land-new}); broadening that coverage is a target for future work
(Section~\ref{sec:future}).

\mapfigure{land_value_existing.png}{Land value per square foot, current assessments}{fig:land-existing}
\mapfigure{land_value_new.png}{Land value per square foot, new assessments}{fig:land-new}

\subsection{Valuation ratio analysis}
To show how a reassessment would redistribute the tax burden, we compute each
parcel's \emph{valuation ratio}---its new assessed value divided by its current
assessed value. Because the county currently assesses at roughly half of market
value, this ratio exceeds one almost everywhere: in aggregate, total residential
assessed value would rise by a factor of about \countyUnderRatio. What matters for
tax fairness is not that common increase but the \emph{dispersion} around
it---tracts rising faster than the typical increase would shoulder a larger share
of the levy. Figure~\ref{fig:valratio} maps the median ratio by census tract,
shaded so the median tract's increase is white; the tracts that would rise most
(deep red) are concentrated in the historically most under-assessed urban core,
while many outlying tracts would rise less than typical.

\mapfigure{valuation_ratio_parcel.png}{Valuation ratio (new $\div$ current assessment) by census tract; white marks the county-median increase}{fig:valratio}

\subsection{Valuation ratio by jurisdiction}\label{sec:jurisdiction}
Aggregating the same valuation ratio to the units that actually levy and govern
property taxes shows who would bear the shift. These published ratios are
normalized so that 1.0 marks the county-typical increase: a jurisdiction above 1.0
would, before any millage adjustment, take on a larger share of the levy, and one
below 1.0 a smaller share. Mapped by Pittsburgh city council district
(Figure~\ref{fig:vr-city}) and county council district (Figure~\ref{fig:vr-county}),
and ranked by school district (Figure~\ref{fig:vr-school}) and municipality
(Figure~\ref{fig:vr-muni}), the pattern is consistent: assessments rise fastest in
the City of Pittsburgh and the older river and mill communities---the most
under-assessed areas today---and rise less than typical across the affluent outer
suburbs. (Municipality and school-district boundaries are not redistributed here,
so those two breakdowns are shown as ranked charts rather than maps.)

\mapfigure{valratio_city_council.png}{Valuation ratio by Pittsburgh city council district (1.0 = county-typical increase)}{fig:vr-city}
\mapfigure{valratio_county_council.png}{Valuation ratio by county council district (1.0 = county-typical increase)}{fig:vr-county}
\figwide{valratio_school_district.png}{width=0.72\linewidth}{Valuation ratio by school district (1.0 = county-typical increase)}{fig:vr-school}
\figwide{valratio_municipality.png}{width=0.80\linewidth}{Valuation ratio by municipality: the 15 highest- and 15 lowest-ranked of 131 (1.0 = county-typical increase)}{fig:vr-muni}

\subsection{Equity analysis}\label{sec:equity}
Because the models use neighborhood income and a market tier as predictors, and
because any reassessment has distributional consequences, we test whether the new
assessments are equitable across neighborhood characteristics. We stratify the
ratio study by Census median-income quintile (Table~\ref{tab:equity-income}) and
by Census-tract racial composition (Table~\ref{tab:equity-minority}).

\begin{table}[htbp]\centering
\caption{Sales ratio and COD by neighborhood income quintile (lowest vs.\ highest).}
\label{tab:equity-income}
\input{generated/equity_income.tex}
\end{table}

\begin{table}[htbp]\centering
\caption{Sales ratio and COD by Census-tract percent-minority band (lowest vs.\ highest).}
\label{tab:equity-minority}
\input{generated/equity_minority.tex}
\end{table}

Two patterns emerge, and they point in different directions. The median sales
ratio is essentially flat across both income quintiles and racial-composition
bands---near 1.0 in every stratum---so the new assessments show no systematic
over- or under-assessment of lower-income or higher-minority neighborhoods. This
is the more legally salient test, and the model passes it. However, assessment
\emph{uniformity} (COD) is consistently worse in the lowest-income and
highest-minority strata. In the most heavily minority tracts the COD exceeds the
IAAO single-family ceiling in all three single-family groups---substantially in
the urban and pre-war groups and more modestly in the suburban group---while the
multi-family group stays within its (looser) standard for small income-producing
residential property. This gradient most likely reflects thinner and more
heterogeneous sales in those neighborhoods rather than bias in the demographic
inputs, but it is a real limitation that any official use of these assessments
must weigh, and a target for future improvement.

\subsection{Effects on tax bills}
A reassessment changes \emph{assessed values}, not directly the dollars owed.
Pennsylvania's anti-windfall provisions require each taxing body---county,
municipality, and school district---to roll its millage rate back after a
countywide reassessment so that total revenue stays essentially flat (the county's
own ordinance requires full revenue-neutrality for county taxes, and state law caps
how much additional revenue school districts and municipalities may capture). In
practice nominal millage rates fall sharply, and a given owner's bill rises or
falls only to the extent that their property's value moved relative to the average
within each taxing body. Most owners, whose relative valuation changes little,
would see little change; the redistribution in Section~\ref{sec:jurisdiction}
operates around a roughly revenue-neutral total. Individual bills also depend on
exemptions and abatements---the Homestead/Farmstead exclusion, LERTA, and senior
and longtime owner-occupant relief among them---that vary across taxing bodies
\citep{acs_tax_exempt}. Because these interactions are genuinely complicated,
Pro-Housing Pittsburgh provides an interactive Assessment Explorer
(\url{https://explorer.prohousingpgh.org}) that lets owners look up our estimated
new value for their parcel. These are preliminary estimates, not official
assessments or tax, legal, or financial advice.

\section{Dataset quality}\label{sec:data-quality}
The models are only as good as the county characteristic data they consume. To
gauge that data's reliability we drew a random sample of \auditSampleN\ residential
parcels and compared key fields against independent references---computed parcel
geometry for lot area, and Zillow listings for the building characteristics
(Table~\ref{tab:audit}). Errors are common but mostly modest. Lot area is the
weakest field, with \auditLotAny\ of \auditSampleN\ parcels off by more than
10\% (we substitute the computed geometry where they disagree); building age,
size, and room counts are usually within one unit or a few percent, and we found no
parcel whose quality or condition grade was blatantly wrong on inspection.
Improving these inputs---lot area and finished living area in particular---is the
most direct route to more accurate assessments.

\begin{table}[htbp]\centering
\caption{Data-quality audit: county values vs.\ independent references (\auditSampleN-parcel sample).}
\label{tab:audit}
\input{generated/audit_table.tex}
\end{table}

\section{Cost estimate}\label{sec:cost}
Reassessment is often assumed to be prohibitively expensive, but the marginal cost
of doing it regularly---with existing data and open-source tools---is modest. Most
of the effort is data cleaning rather than modeling, much of it can be done by less
specialized staff, and once a pipeline is in place each subsequent reassessment
reuses it. Allowing for substantial data-improvement work beyond this preliminary
analysis, we estimate a recurring countywide reassessment can be performed for on
the order of \$1\,million per year. An independent vendor (ValueBase) quoted
roughly \$400{,}000 annually for comparable work, reaching full quality over two to
three years \citep{parkinson2024}.

That is far below the cost of the episodic, court-ordered reassessments of the
past. At a 2024 County Council hearing, Tyler Technologies estimated a one-time
reassessment at over \$30~million on a multi-year timeline \citep{council_tv2024};
at the per-parcel rates discussed there, the county's \totalParcels\ parcels alone
imply tens of millions of dollars. Recent letters to County Council from
public-finance scholars make the same point we do---that with modern methods,
frequent reassessment is far cheaper than the headline figures from past cycles
suggest \citep{anderson2026,skidmore2026,briem2022}. The county should expect to
pay recurring-program prices, not crisis prices.

\section{Limitations}
Several caveats apply. The commercial model is unreliable given sparse sales and
limited characteristic data. The racial-composition breakdown is coarse, because
the percent-minority variable is heavily concentrated near zero in the sales
sample; it should be read as ``highest-minority tracts vs.\ the rest.'' Ratios
are measured against raw sale prices and so partly reflect the 2025--2026 price
trend rather than pure model error. And the land-value decomposition inherits the
land-to-total ratios of the existing (stale) assessments.

Our results also rest on the county's designation of which sales are valid,
arm's-length transactions. An independent University of Chicago analysis of
Allegheny County judged the \emph{current} roll more harshly than we do (reporting
a price-related differential of 1.115 and a COD of 33.11) using different validity
criteria \citep{berry2024}. That divergence suggests the rules for classifying
valid sales deserve their own scrutiny, since they feed directly into both the
models and the ratio studies.

\section{Future work}\label{sec:future}
Priorities include better commercial valuation (additional characteristic and
income data), reducing the uniformity gap in under-traded neighborhoods, a finer
racial-composition equity analysis once enrichment is refined, and per-parcel
prediction intervals to support the appeals process.

\section{Conclusion}
Allegheny County has the data and, now, the tools to end its decades-long
assessment freeze. Using only public data and an open-source toolkit, we produced a
preliminary countywide revaluation that is far closer to market value than the
current roll and, for residential property, reasonably uniform; we quantified how
stale the current assessments have become; we mapped how a reassessment would
redistribute burden across neighborhoods and taxing bodies; we tested the result
for equity across income and racial composition; and we estimated that the program
could be sustained for about \$1\,million per year. The assessments are preliminary
and have real limitations---commercial property, thinly traded neighborhoods, and
the underlying characteristic data chief among them---but they are good enough to
show that regular, transparent, open-source reassessment is both feasible and
affordable. We encourage the County to commit to a regular cycle and to treat this
work as a starting point rather than a finished product.

\appendix
\section{Model details}
This appendix sketches the engines OpenAVMKit combines for each model group. The
aim is to convey the role each plays, not to give a complete treatment.

\paragraph{Multiple regression (MRA).} A linear model estimates sale price as a
weighted sum of property characteristics,
\[ \hat{y}_i = \beta_0 + \beta_1 x_{i1} + \cdots + \beta_k x_{ik}, \]
with coefficients fit by ordinary least squares.
Tables~\ref{tab:coef-suburban}--\ref{tab:coef-multifamily} report the fitted
coefficients for the four residential groups on the dollar scale, so each reads as
the marginal effect of a one-unit change---an added square foot of living area or
land, or a one-step change in a condition or quality grade. (High-cardinality
location terms and the market-value-tier controls are part of the specification but
omitted from the tables for brevity.) The signs are sensible and the spatial lag of
nearby sale prices does much of the work, which makes the linear model a transparent
baseline. In the residential ensembles the linear engines enter in \emph{log} form,
so that blended predictions cannot go negative or compress at the extremes.

\input{generated/coef_tables.tex}

\paragraph{Spatial-lag engines.} Because location dominates value, OpenAVMKit also
forms a distance-weighted average of nearby recent sale prices,
\[ \mathrm{SL}_i = \frac{\sum_{j \in N(i)} w_{ij}\,y_j}{\sum_{j \in N(i)} w_{ij}}, \]
and a variant that first estimates a local price \emph{per unit of area} and scales
it by the subject parcel's size. These capture neighborhood price level using
little or no parcel-specific detail, and serve as strong location benchmarks.

\paragraph{Gradient-boosted trees.} XGBoost \citep{xgboost} and LightGBM
\citep{lightgbm} build an additive ensemble of decision trees, each fit to the
residual left by the trees before it,
\[ \hat{y}_i^{(t)} = \hat{y}_i^{(t-1)} + \eta\,h_t(x_i), \]
with a learning rate $\eta$ and regularization to limit overfitting. Trees capture
nonlinear effects and interactions among characteristics that the linear models
cannot.

\paragraph{Ensemble.} The final value for each parcel is the \emph{median} of a
selected subset of engine predictions; using the median rather than the mean blunts
the influence of any single mis-estimate. OpenAVMKit chooses the subset that
minimizes mean absolute percentage error on held-out sales. For the residential
groups we further restrict the blend to the two log-scale regression engines and
the strongest tree engine: the accuracy-only default tends to retain the regressive
spatial and comparable-sales engines, and dropping them removes most of the
vertical-equity regressivity (Table~\ref{tab:results}) at no cost to dispersion.

\FloatBarrier   % flush any pending figures/tables before the references
\bibliographystyle{plainnat}
\bibliography{references}

\end{document}
